\documentclass[11pt]{article}
\usepackage[a4paper,margin=2.5cm]{geometry}
\usepackage{amsmath,amssymb,booktabs}
\usepackage[hidelinks]{hyperref}
\usepackage{xurl}
\setlength{\parindent}{0pt}
\setlength{\parskip}{4pt}
\newcommand{\MKT}{\mathit{MKT}}
\newcommand{\SMB}{\mathit{SMB}}
\newcommand{\HML}{\mathit{HML}}
\newcommand{\RMW}{\mathit{RMW}}
\newcommand{\CMA}{\mathit{CMA}}
\newcommand{\MOM}{\mathit{MOM}}
\newcommand{\RF}{\mathit{RF}}
\newcommand{\ME}{\mathit{ME}}
\newcommand{\BE}{\mathit{BE}}
\newcommand{\OP}{\mathit{OP}}
\newcommand{\INV}{\mathit{INV}}
\newcommand{\TA}{\mathit{TA}}

\begin{document}
\begin{center}{\large\bfseries Construction of the Fama--French Five Factors}\end{center}

Five factors: $\MKT,\ \SMB,\ \HML,\ \RMW,\ \CMA$ \cite{ff15}. We add $\MOM$ \cite{c97} as a sixth factor in model M4.
Factor returns are downloaded from the Kenneth French Data Library \cite{kf}; the construction below follows \cite{ff15,kf}.

\section*{Abbreviations}
\begin{center}\footnotesize
\begin{tabular}{ll@{\qquad}ll}
\toprule
$\MKT$ & Market excess return & $\ME$ & Market equity (market capitalization) \\
$\SMB$ & Small Minus Big (size) & $\BE$ & Book equity \\
$\HML$ & High Minus Low (value) & $\BE/\ME$ & Book-to-market ratio \\
$\RMW$ & Robust Minus Weak & $\OP$ & Operating profitability \\
$\CMA$ & Conservative Minus Aggressive & $\INV$ & Investment (asset growth) \\
$\MOM$ & Momentum (winners minus losers) & $\TA$ & Total assets \\
$\RF$ & Risk-free rate (1-month T-bill) & FY & Fiscal year \\
$R^m$ & Market return & COGS & Cost of goods sold \\
CRSP & Ctr.\ for Research in Security Prices & SG\&A & Selling, general \& admin.\ expense \\
NYSE & New York Stock Exchange & Rev & Revenue \\
AMEX & American Stock Exchange & $S$ / $B$ & Small / Big (size groups) \\
NASDAQ & NASDAQ Stock Market & $V$ / $N$ / $G$ & Value / Neutral / Growth \\
 & & $R$ / $W$ & Robust / Weak \\
 & & $C$ / $A$ & Conservative / Aggressive \\
 & & $H$ / $L$ & High / Low prior return \\
\bottomrule
\end{tabular}
\end{center}

\section*{Step 1. Stock universe and timing \normalfont\small\cite{kf}}
\begin{itemize}
  \item Universe: CRSP common stocks (share codes 10, 11) listed on NYSE, AMEX, NASDAQ.
  \item Portfolios formed at the end of June of year $t$; held from July $t$ to June $t+1$.
  \item Portfolio returns are value-weighted:
  \begin{equation}
    r_{p,t}=\sum_{i\in p} w_{i,t}\,r_{i,t},\qquad w_{i,t}=\frac{\ME_{i,t-1}}{\sum_{j\in p}\ME_{j,t-1}}
  \end{equation}
\end{itemize}

\section*{Step 2. Sorting variables \normalfont\small\cite{ff15,kf}}
\begin{align}
\ME_i &= \text{price}_i\times\text{shares outstanding}_i && \text{end of June } t\\
(\BE/\ME)_i &= \frac{\BE_{i,\,\text{FY }t-1}}{\ME_{i,\,\text{Dec }t-1}}\\
\OP_i &= \frac{\text{Rev}-\text{COGS}-\text{Interest expense}-\text{SG\&A}}{\BE+\text{Minority interest}} && \text{FY } t-1\\
\INV_i &= \frac{\TA_{i,t-1}-\TA_{i,t-2}}{\TA_{i,t-2}}
\end{align}

\section*{Step 3. Breakpoints and $2\times3$ portfolios \normalfont\small\cite{ff15,kf}}
Size splits at the NYSE median $\ME$. Each of $\BE/\ME$, $\OP$, $\INV$ splits at its NYSE 30th and 70th percentiles.
Intersecting gives three independent sets of six portfolios.

\begin{center}
\begin{tabular}{lccc}
\toprule
Sort & Top 30\% & Middle 40\% & Bottom 30\% \\
\midrule
$\BE/\ME$ & Value ($V$) & Neutral ($N$) & Growth ($G$) \\
$\OP$ & Robust ($R$) & Neutral ($N$) & Weak ($W$) \\
$\INV$ & Aggressive ($A$) & Neutral ($N$) & Conservative ($C$) \\
\bottomrule
\end{tabular}
\end{center}
Each portfolio is labeled Small ($S$) or Big ($B$) plus its group, e.g.\ $SV$ = small value.
Each label in the formulas below denotes that portfolio's value-weighted return $r_{p,t}$.

\section*{Step 4. The five factors}

\textbf{(1) Market} \cite{ff93}
\begin{equation}
\MKT_t = R^m_t-\RF_t
\end{equation}
$R^m_t$: value-weighted return of all stocks in the universe;\ $\RF_t$: one-month Treasury bill return.

\textbf{(2) Size: Small Minus Big} \cite{ff15}
\begin{align}
\SMB_{B/M} &= \tfrac13(SV+SN+SG)-\tfrac13(BV+BN+BG)\\
\SMB_{OP}  &= \tfrac13(SR+SN+SW)-\tfrac13(BR+BN+BW)\\
\SMB_{INV} &= \tfrac13(SC+SN+SA)-\tfrac13(BC+BN+BA)\\
\SMB &= \tfrac13\left(\SMB_{B/M}+\SMB_{OP}+\SMB_{INV}\right)
\end{align}

\textbf{(3) Value: High Minus Low book-to-market} \cite{ff15}
\begin{equation}
\HML = \tfrac12(SV+BV)-\tfrac12(SG+BG)
\end{equation}

\textbf{(4) Profitability: Robust Minus Weak} \cite{ff15}
\begin{equation}
\RMW = \tfrac12(SR+BR)-\tfrac12(SW+BW)
\end{equation}

\textbf{(5) Investment: Conservative Minus Aggressive} \cite{ff15}
\begin{equation}
\CMA = \tfrac12(SC+BC)-\tfrac12(SA+BA)
\end{equation}

\section*{Additional factor: momentum \normalfont\small\cite{c97,kfmom}}
Formed monthly at the end of $t-1$; size split at NYSE median, prior-return split at NYSE 30th/70th percentiles.
\begin{align}
r^{\text{prior}}_{i,t} &= \prod_{s=t-12}^{t-2}(1+r_{i,s})-1 && \text{skips month } t-1\\
\MOM_t &= \tfrac12(SH+BH)-\tfrac12(SL+BL) && H/L:\ \text{high/low prior return}
\end{align}

\begin{thebibliography}{9}\small
\bibitem{c97} Carhart, M. M. (1997). On persistence in mutual fund performance. \emph{Journal of Finance}, 52(1), 57--82.
\bibitem{ff93} Fama, E. F., \& French, K. R. (1993). Common risk factors in the returns on stocks and bonds. \emph{Journal of Financial Economics}, 33(1), 3--56.
\bibitem{ff15} Fama, E. F., \& French, K. R. (2015). A five-factor asset pricing model. \emph{Journal of Financial Economics}, 116(1), 1--22.
\bibitem{kf} French, K. R. Data Library: Fama/French 5 Factors (2x3); Variable Definitions.
\url{https://mba.tuck.dartmouth.edu/pages/faculty/ken.french/Data_Library/f-f_5_factors_2x3.html}
\bibitem{kfmom} French, K. R. Data Library: Momentum Factor (Mom).
\url{https://mba.tuck.dartmouth.edu/pages/faculty/ken.french/Data_Library/det_mom_factor.html}
\end{thebibliography}
\end{document}
